\chapter{Журнал опытов: SI, PDG и реестр сверок}
\label{app:si-pdg-experiments}

\section{Зачем это приложение}
\label{sec:si-pdg-experiments}

Глава~\ref{ch:si-sm} фиксирует \emph{формулы} вывода ($v$, массы, $\sin^2\theta_W$, $\alpha_s$).
Сходство с СИ и таблицами PDG/CODATA --- не постулат и не «мы решили, что всё ок'',
а \textbf{серия вычислительных экспериментов} серии CE-M
(\cref{ch:computational-laboratory,ch:critical-construction,sec:lab-series-ce-m}).

Каждая строка ниже --- косвенное или алгебраическое сравнение с эталоном;
методика $u$, $E_n$, анти-круговость --- \cref{ch:measurement-processing}.
Аннигиляция и этажи --- серии CE-A / CE-F (\cref{ch:critical-floors}).

\section{Таблица CE-M (автогенерация)}
\label{sec:si-pdg-table}

% Auto-generated by scripts/gen_lab_records_tex.py — do not edit.
% Generated: 2026-09-29 20:24 UTC
% Journal prose: book/sources/lab-journal/ (hand-authored).
\begin{table}[htbp]
\centering
\caption{Серия CE-M: модель (upstream §8.2) vs эталон PDG/CODATA. Не натурный эксперимент --- воспроизводимый прогон реестра.}
\label{tab:si-pdg-ce-m}
\small
\begin{tabular}{@{}lrrrl@{}}
\toprule
Наблюдаемая & Модель & Эталон & $|\delta|$ & Verify \\
\midrule
$v$ [GeV] & 246{,}09 & 246{,}22 & $5{,}39\times 10^{-4}$ & \texttt{Higgs\_mass} \\
$m_H$ [GeV] & 125{,}31 & 125{,}25 & $4{,}73\times 10^{-4}$ & \texttt{Higgs\_mass} \\
$m_p$ [GeV] & 0{,}9353 & 0{,}9383 & $3{,}16\times 10^{-3}$ & \texttt{Proton\_mass} \\
$m_e$ [GeV] & 0{,}000513 & 0{,}000511 & $4{,}50\times 10^{-3}$ & \texttt{Electron\_mass} \\
$m_n$ [GeV] & 0{,}9363 & 0{,}9396 & $3{,}44\times 10^{-3}$ & \texttt{Neutron\_mass} \\
$\alpha^{-1}$ & 137{,}035999 & 137{,}035999 & $7{,}19\times 10^{-11}$ & \texttt{Alpha\_si\_bridge} \\
$\sin^2\theta_W$ (bare) & 0{,}23077 & 0{,}23122 & $1{,}95\times 10^{-3}$ & \texttt{Alpha\_bridges} \\
$\alpha_s(M_Z)$ & 0{,}11788 & 0{,}11790 & $1{,}36\times 10^{-4}$ & \texttt{---} \\
\bottomrule
\end{tabular}
\end{table}


\section{Воспроизведение}
\begin{itemize}
  \item Обновить числа в книге:
    \texttt{python scripts/gen\_lab\_records\_tex.py} из корня репозитория.
  \item Прогнать реестр:
    \texttt{python verify\_principles.py --suite alpha --suite sm --device cpu}
    (полный ship: \texttt{--ship}).
  \item Две скорости $c_0$ и $c$ --- CE-M-00 (\cref{lab:ce-m-00,rem:two-c-burden}); макро-фронт --- CE-M-01.
  \item Полные постановки CE-M-00--03 --- в конце тома \emph{construction},
    глава «Критические эксперименты''; поля \textbf{Анализ} подставляются из тех же прогонов.
\end{itemize}
